\documentclass[AutoFakeBold]{ctexart}	% AutoFakeBold 伪黑体加粗

%*****************超链接***********************************
% colorlinks 将超链接以颜色来标识，默认用方框来标识
% url颜色、引用颜色
\usepackage[colorlinks,urlcolor=black,linkcolor=black]{hyperref}

%*****************编号**********************************
\usepackage{enumerate}
\usepackage{enumitem}
\setlist[enumerate,1]{label=(\arabic*),labelindent=\parindent,leftmargin=*}	% 表示对一级列表进行设置(1)
\setlist[enumerate,2]{label=(\alph*)}	% 表示对一级列表进行设置(a)

%*****************数学公式**********************************
\usepackage{amssymb}
\usepackage{amsmath}

%*****************图片及其相关设置***************************
\usepackage{graphicx}
\graphicspath{{figures/}}
\usepackage{subfigure}	% 子图

%*****************表格及其相关设置***************************
\usepackage{booktabs}	% 三线表
\usepackage{multirow}
\usepackage{array}

%*****************代码环境************************************
\usepackage{listings}
\usepackage{xcolor}
\lstset{
	numbers=left, %设置行号位置
    numberstyle=\tiny, %设置行号大小
    keywordstyle=\color{blue}, %设置关键字颜色
    commentstyle=\color[cmyk]{1,0,1,0}, %设置注释颜色
    frame=single, %设置边框格式
    escapeinside=``, %逃逸字符(1左面的键)，用于显示中文
    breaklines, %自动折行
    extendedchars=false, %解决代码跨页时，章节标题，页眉等汉字不显示的问题
    xleftmargin=2em,xrightmargin=2em, aboveskip=1em, %设置边距
    tabsize=4, %设置tab空格数
    showspaces=false %不显示空格
}

%*****************题注及其相关设置***************************
\usepackage{caption}
\DeclareCaptionFont{heiti}{\heiti}	% 声明题注的字体
% labelsep=space 将'表1: XXX' 设置为 '表1 XXX '
% [small, bf, heiti] 黑体 小5 加粗, heiti 前面使用前面定义的 \heiti
\captionsetup{labelsep=space,font={small,bf,heiti},skip={5pt}}	

%*****************页边距、行间距、字体**********************
\usepackage[left=2.5cm,right=1.97cm,top=2.5cm,bottom=2.5cm]{geometry}	% 页边距
\renewcommand{\baselinestretch}{1.25}	% 行间距

%*****************全局英文字体设置为 Times New Roman *******
\usepackage{fontspec}
\setmainfont{Times New Roman}

%*****************多级标题格式修改**************************
\ctexset{
	section={format={\heiti \zihao{4} \bfseries}},	% 加粗
	subsection={format={\heiti \zihao{5} \bfseries},beforeskip=0pt,afterskip=0pt},	% 段前0pt、段后0pt
	subsubsection={format={\kaishu \zihao{5}},beforeskip=0pt,afterskip=0pt}	% 段前0pt、段后0pt
}

%%*****************标题************************************
%\title{
%    \heiti \zihao{2}用\LaTeX 排版 
%    \footnotetext{	% 脚标
%        \hskip -6pt 投稿日期：YYYY-MM-DD \\	% 向前移动 6pt
%        \hspace*{1em}项目名称:	% 向后移动 1em
%        \begin{minipage}[t]{0.85 \linewidth}	% [t]顶部对齐, 0.85 的行宽
%            国家自然科学基金资助项目1，国家自然科学基金资助项目2，国家自然科学基金资助项目3，国家自然科学基金资助项目4，国家自然科学基金资助项目5，国家自然科学基金资助项目6，国家自然科学基金资助项目7
%        \end{minipage}}}
%% 作者、联系方式
%\author{
%        \kaishu \zihao{-4} 作者 \\	% 楷体 小四
%        \songti \zihao{-5} 学校\qquad 地址
%}
%\date{}	% 不显示日期
%%%%%%%%%%%%%%%%%%%%%%%%%%%%%%%%%%%%%%%%%%%%%%%%%%%%%%%%%%%%%%%%%%%%%%%
\begin{document}
	\maketitle	% 显示标题
	\vspace{-10pt}%---------------------摘要------------------------
%\noindent 取消首行缩进
% 中文摘要、关键字、分类号
\noindent \heiti 摘要：\songti 本文主要针对\LaTeX 为中文论文排版进行简要的描述，以工业工程该期刊的要求作为模板。 \\
\heiti 关键词：\songti \LaTeX ； 中文期刊 ；工业工程 \\
\heiti 中图分类号：\songti 000 \hspace{3cm}     \heiti 文献标识码：\songti A \\

% 英文标题居中加粗、作者、简介
\begin{center}
	\textbf{\zihao{4} How to Write An Chinese Journal with \LaTeX} \\ 
	\zihao{5} Joefsong \\ 
	\zihao{-5} (Heiheihei College Papapa School, QQ Group 970479548) \\
\end{center}

% 英文摘要和关键字
\noindent \textbf{Abstract:} Wo de ying yu hen lan, jiu bu fan yi le. Fan zheng wo jue de pai ban he fan yi ye mei you shen me guan xi. Dan shi wei le rang ying wen zhai yao nei rong xian de feng man hai shi yao duo xie yi dian. hao le wo bian bu xia qu le. Jiu zhe yang ba.\\
\textbf{key words:}  \LaTeX ; Chinese Journal; Industrial Engineering Journal \\

% 引言
现存的用来编辑或排版的工具很多，最被熟知的应该是微软的软件，笔者了解到国内商用(编辑部)比较多的应该是方正书版。但是微软的软件实现很多功能都需要插件，例如：word中插入公式所用到的内置公式编辑器较为繁琐，而Mathtype插件的兼容性很差，且需要收费才能使用正版插件，此外在处理文本的过程中经常出现格式问题，如无法与文本对齐等，而且当部分公式的字体(如表格内公式)需要调整字号的时候需要进入公式编辑器逐个调整；除此之外，在插入参考文献的时候需要借助Endnote 等工具；在作图的时候需要借助Viso等。而方正书版最大的缺陷就是收费，且价格不菲。因此完全免费且容易上手的排版软件\LaTeX 脱颖而出，并且英文期刊几乎全部都要求用\LaTeX 排版，部分中文期刊也都接受\LaTeX 排版的文章，少部分中文期刊要求用\LaTeX 排版，也有一部分中文期刊仅接受word的排版。因为工业工程期刊是接受\LaTeX 排版的论文的，并且提供了\LaTeX 的模板，但是由于版本问题，无法兼容现在的texlive。因此本文选择了工业工程这个期刊为例，按照其提供的论文排版要求，对该篇论文进行人工的排版。最终为需要利用\LaTeX 排版中文期刊论文的人士提供有效易学的方法。	% 标题与摘要之前的间隔减小 10pt
	%---------------------正文------------------------
%%%%%%%%%%%%%%%%%%%%%%%%%%%%%%%%%%%%%%%%%%%%%%%%%%%%%%%%%%%%%%%%%%%
\section{基本文本格式以及段落调整}	% 一级标题
期刊论文的要求中，要求最多的就是字体字号以及章节格式，因此本节简要介绍\LaTeX 中的中英文字体字号设置以及段落格式的基本要求的实现。

\subsection{字体和字号}	% 二级标题
\LaTeX 中若使用了ctex或者ctexart宏包，则需要注意\LaTeX 包含常用的四种中文字体以及三种英文字体可供选择。由于期刊论文一般不要求其他字体，因此本节不在讲述载入系统字体的方法。

\subsubsection{字体}%三级标题
\LaTeX 默认的中文字体为宋体， 且有黑体、 楷书和仿宋三种气体可供选择； 默认的英文字体为罗马字体 (Roman Family) ，且有无衬线字体 (Sans Serif Family) 以及打印字体 (Typewtiter Family) 可供选择。具体字体及其命令如下：

\begin{table}[htbp] \zihao{-5} 

	\caption{字体命令} \label{字体命令}
	\centering
	\begin{tabular}{l p{3cm}<{\raggedleft} l p{3cm}<{\centering}}
		\toprule
		中文字体&命令& 英文字体& 命令\\
		\hline
		宋体&	songti&	罗马字体& rmfamily\\
		黑体	& heiti &	无衬线字体 &	sffamily\\
		楷体 &	kaishu &	打印字体& 	ttfamily\\
		仿宋& fangsong & \\
		\hline
	\end{tabular}
\end{table}
	
\subsubsection{字号}
由于使用了中文的宏包，因此在设置字号的时候可以直接使用 zihao 命令进行定义字号大小。大括号中是字号的大小，如五号字体为5，小五号字体为-5。


\subsection{段落调整}
\subsubsection{基本段落格式}
需要注意的是\LaTeX 在编辑中文文本的时候会忽略空格和回车，两个中文字符之间不管有多少空格，生成的PDF文件都不会显示出来。此外单个回车也不会是生成全新段落，若要实现生成新的段落需要在段落间空一行。若仅需要换行而不需要生成新的段落可以用双反斜杠实现。利用空行生成的新段落会默认首行缩进两个字符，利用双反斜杠生成的新段落不会首行缩进。若想要取消某段落首行缩进，须在段首使用 noident 命令。而中文文本和英文文本在一起出现的 时候，会自动在两者之间生成空格。quad 一个空格  qquad 两个空格  

\subsubsection{行间距调整}
\LaTeX 使用中文ctex 宏包之后，默认的是1.3倍行距，而中文期刊不会对行距有特别的要求，因此可直接使用默认行距，若需要设置论文行距，可以在导言区重新定义baselinestretch的值，具体参见导言区命令。各级标题与段落之间的间距设置将在\ref{多级标题}%讲多级标题的章节
给出。

也会遇到情况如摘要和作者的间距太长，仅需提高摘要的位置，使其靠近作者或者标题，此实仅需利用vspace命令定义向上或向下移动的距离。常见距离单位如下：也可通过baselineskip 或 textwidth 等定义相对距离。
单位 	名称 	说明
mm 	毫米 	1 mm = 2.845 pt
pt 	点 	1 pt = 0.351 mm
cm 	厘米 	1 cm= 10 mm= 28.453 pt
in 	英寸 	1 in = 25.4 mm = 72.27 pt
ex 	ex 	1 ex = 当前字体尺寸中 x 的高度
em 	em 	1 em = 当前字体尺寸中 M 的宽度

\subsubsection{自动编号} \label{自动编号}
若在论文写作过程中有分条的需求可利用itemize或者enumerate 环境实现分条和交叉引用。itemize 可以满足自定义的序号标识，而enumerate则可以自动生成有序的标号。类似的\LaTeX 中也包含有定理定义证明等环境可以自动生成序号以及可以实现交叉引用，但是由于在期刊论文中的实用性不高，因此本文不予详细介绍。

\section{多级标题以及文件导入} \label{多级标题}
\LaTeX 在编辑过程中，通过命令可以自动生成有序的标题，并可以在导言区对其格式进行设置。然而在写文章的过程中，没有人可以一气呵成，然而由于在利用\LaTeX 写作过程中，会出现很多命令，并且当篇幅过长的时候影响审阅，以及会使运行时间变得越来越长，因此将不同章节份文件保存，最后进行整合可以使文章看起来简介有序，并且在无需更改的时候可以先行注释不予运行，加快运行速度，提高编辑效率。 
\subsection{多级标题格式设置}
\LaTeX 默认有三级标题，而在实际论文的写作过程中很少用到四级标题，因为会使标题序号变得异常的繁琐且不美观，如果需要四级标题，可直接利用\ref{自动编号}%讲自动编号的章节
中的itemize 环境实现。

而三级标题在未使用ctex中文宏包之前，默认为左对齐。在使用ctex宏包之后依然为左对齐，但是在使用ctexart宏包之后会默认为居中对齐。由于本文使用的文件类型为ctexart因此在需要使一级标题左对齐的时候需要在导言区进行设置。各级标题及其对应名命令如下：\\
一级标题 section\\
二级标题 subsection\\
三级标题 subsubsection

对各级标题格式进行设置可以利用ctexsetup语句，但是此语句并非实现此功能的唯一方法，其他方法不予详细介绍。本文主要针对工业工程期刊的要求，对各级标题进行字体字号设置，此外针对标题与段落间距的设置也会进行简要介绍，由于笔者暂未发现其他设置需求，因此本文不予介绍。
\subsection{文件导入}
实现章节以单独文件的形式保存，仅需将章节及文本内容单独一tex类型文件保存，且在正文中利用input 命令或者include 命令进行载入。 需要注意的是，当章节内容单独保存的时候不需要(不能)添加导言区内容；此外，利用input 命令载入的时候，文本会直接在段落后载入，而include 载入的时候则会在新的页面生成段落章节。因此可以根据自己需要选择使用，但期刊论文一般使用input命令即可。
	\section{插入公式}
	\LaTeX 相比于word 的优点之一便是公式的兼容性，本文仅对数学符号以及相关公式进行简要的介绍。
	\subsection{数学符号}
	在使用数学符号之前需要在导言区使用amssymb宏包，相关命令不予赘述，可以在texstudio 的符号中直接进行选择插入，无需记忆。需要注意的是，在使用数学符号的时候需要在 \$ \$ 内输入。 且由于反斜杠和 dollar 符号以及大括号等特殊符号与 \LaTeX 命令符号存在冲突，这些符号是 \LaTeX 的保留字符，因此在需要输入这些符号的情况下，有特殊的方式，可以参考下面三个实力，更多的保留字符可以自行搜索：\\
%	\hspace*{3cm}
	dollar 符号 ~~ \$  \hspace{5cm}	反斜杠 ~~ $\backslash$ \hspace{5cm}大括号    $\left\lbrace \right\rbrace$	 
	
	需要注意的是\$ \$ 内的符号都是默认的斜体，因此如需使用正体符号，需要重新定义。
	\subsection{常见公式} 
	常见的公式种类分为行内公式和行间公式，行内公式可以直接使用 \$ \$符号进行书写；而行间公式如$$a+b=c$$可通过\$\$ \quad \$\$符号实现，如果需要对公式进行标号以及交叉引用，可以通过equation环境实现。如果需要连续编辑多行公式则可以借助gather环境来实现。示例如下：
多行公式：
	\begin{gather}
		a+b=c \\
		\alpha + \beta =\gamma \notag \\
		x+y=z \notag
	\end{gather}
	方程组或者分段函数可以用cases环境实现:
	\begin{equation}
		\begin{cases}
			a+b=c \\
			\alpha + \beta =\gamma  \\
			x+y=z
		\end{cases} 
	\end{equation}
	\begin{equation}
		f(x)= 
		\begin{cases}
			ax+b+c+1 &if x<\gamma \\
			\alpha x + \beta &if x>\gamma 
		\end{cases} 
	\end{equation}
	长公式用split环境实现
	\begin{equation}
		\begin{split}
			f(x)&=10000000a^2x+123456789bx+1010101010c+1010101010+\alpha x + \beta + \gamma\\
			&=15746987\sqrt{\alpha x + \beta}+121212121212\lambda^{ax}x+\frac{1}{\lambda^2}\\
			&=0
		\end{split}
	\end{equation}
	
	由于矩阵的用法并不复杂，因此不再单独介绍，如需要了解，附录中有网友对矩阵输入整理的链接。
	\section{插入图片}
期刊论文很少有纯文字的叙述，很多情况下都要有图像加以证明。而\LaTeX 中插入图片也非常的方便，且接受tif 或 eps矢量图的格式。下面就介绍几种常见插入图片的格式。

\subsection{一般图片插入}
本节介绍单张图片的插入以及相关题注和交叉引用，在插入图片前需在导言区使用graphicx宏包。如图\ref{头像}
所示为单张图像的插入。需要注意的是需要插入的图片需要和tex 文件在同一个文件夹，为了方便管理也可指定文件夹，并把文章所需要的图片都放在相同文件夹内便于管理;需要插入的图片文件名不能存在中文字符。如果想实现图片包含题注以及可以交叉引用，则需要使用figure环境。

\begin{figure}[htbp]
	\centering
	\includegraphics[width=0.4\linewidth]{image}	% 图片的宽度为 0.4 倍的行宽
	\caption{头像}
	\label{头像}
\end{figure}

\subsection{含有子图的图片}
在期刊论文写作的过程中，可能单张图片无法满足想要达到的效果。更多情况下我们需要对多张图片进行综合分析或对比分析，因此本节主要分享三种包含子图的插入方式。
\subsubsection{多张图片一个题注}
本节介绍多张图片共用同一张标题，比如我用过很多张头像，而且我并不需要对头像进行分类说明，只需要概括说明，此实可以直接插入多张图片如图\ref{用过的头像}，此时仅有一个题注。两个图片的位置距离都遵循段落格式的基本设置。

\begin{figure}[htbp]
	\centering
	\includegraphics[width=0.4\linewidth]{image} \hspace{1cm}	% 图片间隔 1cm
	\includegraphics[width=0.4\linewidth]{gou}
	\caption{用过的头像}
	\label{用过的头像}
\end{figure}
\subsubsection{多张图片拥有小题注以及公共题注}
上节所表述的方法是针对于不需要区分子图情况，因此当需要区分子图的时候，例如引用某一个子图的时候，就需要在拥有公共题注的时候对每个子图进行命名，比如图\ref{情侣头像}
为情侣头像，图\ref{情侣狗}
为生肖为狗的头像，图\ref{生肖鼠}
为生效为鼠的头像。
\begin{figure}[htbp]
	\centering
	\subfigure[情侣狗]{\label{情侣狗}\includegraphics[width=0.4\linewidth]{gou}} \quad
	\subfigure[情侣鼠]{\label{生肖鼠}\includegraphics[width=0.4\linewidth]{shu}}
	\caption{情侣头像}
	\label{情侣头像}
\end{figure}
\subsubsection{多张图片并排显示拥有单独题注}
有时候并列图片并非是用来做综合分析或者差异分析的，也有可能只是为了减单个图片所占用的空间，因此需要是多个图片并排显示，并且分别有各自的题注。由于直接用figure环境，其题注都是居中显示且单独成段，因此需要借助minipage 的环境插入图片。 例如图\ref{生肖虎}
和图\ref{生肖鸡}
为并列显示的两张拥有各自题注的图片。
\begin{figure}[htbp]
	\centering
	\begin{minipage}[htbp]{0.48\linewidth}
		\centering
		\includegraphics[width=0.8\linewidth]{hu}	% 0.8 是 minipage 的行宽
		\caption{生肖虎} \label{生肖虎}
	\end{minipage}
\begin{minipage}[htbp]{0.48\linewidth}
	\centering
	\includegraphics[width=0.8\linewidth]{ji}
	\caption{生肖鸡} \label{生肖鸡}
\end{minipage}
\end{figure}
	%--------------------------------附录-----------------------------------
附录：
\begin{enumerate}
\item 图片：\\		
	\url{https://www.cnblogs.com/yifdu25/p/8326025.html}\\		
	\url{http://www.ctex.org/documents/latex/graphics/graphics.html}
\item 题注 caption 格式：\\
	\url{https://www.latexstudio.net/archives/8652.html}
\item Excel2latex 宏：\\
	\url{https://www.ctan.org/tex-archive/support/excel2latex/}
\item 特殊符号：\\
	win+R 键入 texdoc symbols
\item 数学公式和矩阵:\\
	\url{https://www.jianshu.com/p/fda4c06d43cc}\\
	\url{https://www.cnblogs.com/solvit/p/11345482.html#array\%E7\%8E\%AF\%E5\%A2\%83}
\item 脚注的使用方法以及如何修改格式：\\
	\url{https://www.latexstudio.net/archives/51620.html}\\
	\url{https://blog.csdn.net/jorg_zhao/article/details/80703981}\\
	\url{https://www.latexstudio.net/hulatex/package/note.htm}
\item 字体加粗和新罗马字体设置:\\
	\url{https://www.zhihu.com/question/58456658/answer/157120688}\\
	\url{https://tex.stackexchange.com/questions/79935/how-to-change-the-spacing-in-latexit-for-times-new-roman-font}
\item 导入章节input 以及 include 的用法及区别：\\		
	\url{https://blog.csdn.net/weixin_42919606/article/details/82939495}	
\item  标题格式更改\\		
	\url{https://www.jianshu.com/p/d7848f815e5f}		
\item 编号方法itemize 和 enumerate ：\\		
	\url{http://blog.sina.com.cn/s/blog_77f5a65c0101fmjl.html}\\		
	\url{https://blog.csdn.net/qq_18055167/article/details/83714725}		
\item \label{字体字号} 字体字号及相关设置：\\		
	\url{https://blog.csdn.net/weixin_44537194/article/details/87720878}\\	
	\url{https://blog.csdn.net/ujsdui/article/details/79075327}
\item 常用命令：\\		
	\url{https://blog.csdn.net/garfielder007/article/details/51646802}
\item 文档类型以及页面设置的基本介绍：\\
	\url{https://blog.csdn.net/wei_love_2017/article/details/86617235}
\end{enumerate}		
\end{document}